\section{Графы и их представление в ЭВМ}
\label{sec:graphs-representation}

\subsection{Основные определения}

Графом называется пара
\[
    G=(V,E),
\]
где $V$ --- множество вершин, а $E$ --- множество рёбер. В неориентированном
графе ребро является неупорядоченной парой $\{u,v\}$, в ориентированном графе
дуга является упорядоченной парой $(u,v)$.

Граф может быть простым, мультиграфом, псевдографом, ориентированным,
взвешенным или помеченным. В простом графе нет петель и кратных рёбер.
Число вершин обозначим $n=|V|$, число рёбер --- $m=|E|$.

Вершины $u$ и $v$ называются смежными, если они соединены ребром. Степень
вершины $v$ --- число инцидентных ей рёбер:
\[
    d(v)=|\{e\in E: v\in e\}|.
\]
Для орграфа различают полустепень захода $d^-(v)$ и полустепень исхода
$d^+(v)$.

Для неориентированного графа справедлива лемма о рукопожатиях:
\[
    \sum_{v\in V}d(v)=2m.
\]
Для орграфа
\[
    \sum_{v\in V}d^+(v)=\sum_{v\in V}d^-(v)=m.
\]

\subsection{Маршруты, цепи и циклы}

Маршрутом называется чередующаяся последовательность вершин и рёбер, в которой
каждое ребро соединяет соседние вершины последовательности. Если рёбра не
повторяются, маршрут называется цепью; если не повторяются вершины, кроме
возможного совпадения первой и последней, --- простой цепью.

Замкнутая простая цепь называется циклом. Граф без циклов называется
ациклическим. Две вершины связаны, если существует соединяющая их цепь.
Связный неориентированный граф имеет одну компоненту связности.

Расстояние $d(u,v)$ --- длина кратчайшей цепи между вершинами. Эксцентриситет,
радиус и диаметр графа определяются формулами
\[
 e(v)=\max_{u\in V}d(u,v),\qquad
 R(G)=\min_{v\in V}e(v),\qquad
 D(G)=\max_{v\in V}e(v).
\]

\subsection{Матрица смежности}

Для графа с вершинами $v_1,\ldots,v_n$ матрица смежности $A=(a_{ij})$ задаётся
как
\[
 a_{ij}=\begin{cases}
 1, & \text{если существует ребро или дуга }(v_i,v_j),\\
 0, & \text{иначе}.
 \end{cases}
\]
Для простого неориентированного графа $A$ симметрична, а диагональные элементы
равны нулю. Для взвешенного графа вместо единицы обычно хранят вес ребра.

Память: $O(n^2)$. Проверка наличия ребра выполняется за $O(1)$, но перебор
соседей вершины требует $O(n)$. Представление выгодно для плотных графов.
Степени матрицы имеют комбинаторный смысл: элемент $(A^k)_{ij}$ равен числу
маршрутов длины $k$ из $v_i$ в $v_j$.

\subsection{Матрица инцидентности}

Матрица инцидентности $B$ имеет размер $n\times m$. Для неориентированного графа
\[
 b_{ij}=\begin{cases}
 1, & \text{если вершина }v_i\text{ инцидентна ребру }e_j,\\
 0, & \text{иначе}.
 \end{cases}
\]
Каждый столбец простого графа содержит две единицы. Для орграфа часто используют
знаковую матрицу: $-1$ в начале дуги и $+1$ в конце. Память в плотном виде
составляет $O(nm)$, поэтому на практике матрица инцидентности чаще хранится
разреженно.

\subsection{Списки смежности}

Для каждой вершины хранится список её соседей. Память составляет
\[
    O(n+m)
\]
для орграфа и $O(n+2m)=O(n+m)$ для неориентированного графа.
Перебор всех соседей вершины выполняется за $O(d(v))$, а полный обход графа ---
за $O(n+m)$. Проверка конкретного ребра в обычном списке занимает $O(d(v))$;
при использовании хеш-таблиц ожидаемое время можно уменьшить до $O(1)$.

Списки смежности являются стандартным представлением разреженных графов и
используются в BFS, DFS, алгоритмах Дейкстры, Краскала, Прима и других.

\subsection{Список рёбер}

Граф представляется массивом пар или троек
\[
    (u_i,v_i)\quad\text{или}\quad(u_i,v_i,w_i),\qquad i=1,\ldots,m.
\]
Память --- $O(m)$. Представление удобно, когда алгоритм обрабатывает рёбра
последовательно, например в алгоритме Краскала или Беллмана--Форда. Поиск всех
соседей одной вершины без дополнительного индекса требует просмотра всего
массива.

\subsection{Выбор представления}

\begin{center}
\renewcommand{\arraystretch}{1.25}
\begin{tabular}{|l|c|c|c|}
\hline
Представление & Память & Проверка ребра & Перебор соседей \\
\hline
Матрица смежности & $O(n^2)$ & $O(1)$ & $O(n)$ \\
Списки смежности & $O(n+m)$ & $O(d(v))$ & $O(d(v))$ \\
Список рёбер & $O(m)$ & $O(m)$ & $O(m)$ \\
\hline
\end{tabular}
\end{center}

Для плотного графа, когда $m$ порядка $n^2$, удобна матрица смежности. Для
разреженного графа предпочтительны списки смежности. В реальных программах
часто применяют комбинированные структуры: массив вершин, массив рёбер и
индексы начала списков смежности.

\subsection{Изоморфизм графов}

Графы $G_1=(V_1,E_1)$ и $G_2=(V_2,E_2)$ изоморфны, если существует биекция
$\varphi\colon V_1\to V_2$, сохраняющая смежность:
\[
    \{u,v\}\in E_1
    \quad\Longleftrightarrow\quad
    \{\varphi(u),\varphi(v)\}\in E_2.
\]
Изоморфные графы имеют одинаковые инварианты: числа вершин и рёбер, набор
степеней, число компонент, циклов и так далее. Обратное в общем случае неверно.

\subsection{Что важно сказать на экзамене}

Следует дать определения графа, степени, пути и цикла, затем сравнить четыре
основных представления. Главный практический выбор: матрица смежности для
плотных графов и быстрых запросов о наличии ребра; списки смежности --- для
разреженных графов и линейных обходов.
